\documentclass[10pt,a4paper]{amsart}
\usepackage[margin=19mm]{geometry}
\usepackage{amsmath,amssymb,mathtools}
\usepackage[scr=rsfs,cal=euler]{mathalfa}
\usepackage{xfrac}
\usepackage[unicode,hidelinks,bookmarks=false]{hyperref}
\newcommand{\ie}{i.e.\ }
\newcommand{\cf}{cf.\ }
\newcommand{\resp}{respectively\ }
\newcommand{\Subsection}{\subsection}
\providecommand{\nobreakdash}{\nobreak\hbox{-}}
\setlength{\emergencystretch}{2em}
\multlinegap0pt
\usepackage{amssymb}
\usepackage{tikz}
\usepackage{bm}
\usepackage{algorithm}\usepackage{algpseudocode}
\usepackage{enumerate}
\newtheorem{lemma}{Lemma}
\newtheorem{theorem}{Theorem}
\title{Extremely Amenable Automorphism Groups of Countable Structures}
\author{Mahmood Etedadialiabadi, Su Gao, Feng Li, Ruiwen Li}
\date{Source: arXiv:2411.02760v2 (CC BY 4.0)}
\begin{document}
\maketitle

\noindent\emph{Source and licence.} Extremely Amenable Automorphism Groups of Countable Structures. Version arXiv:2411.02760v2 (CC BY 4.0), distributed by arXiv under the Creative Commons Attribution 4.0 International licence, http://creativecommons.org/licenses/by/4.0/, as recorded in the arXiv licence metadata for this exact version and shown on the abstract page https://arxiv.org/abs/2411.02760v2. Full licence notice: \texttt{licenses/arxiv-2411.02760-CC-BY-4.0.txt}.\par\smallskip

\noindent\textit{Editorial scope.} Counted unit: the article's theorem thm:CDVSinfty (source0.tex lines 936--937). The article's complete original proof of it is delivered with it (source0.tex lines 939--982, one \texttt{proof} environment of 44 lines). No mathematics is written, repaired or re-derived: the statement and the proof are the article's own lines, in the article's own notation, and no retained line is substituted or re-flowed.  Retained uncounted as the source-local context the proof needs: lines 891--892.  Nothing was omitted from any retained span, so the package carries no omission record.  No retained line contains a citation, so the wrapper carries no bibliography and no bibliography entry.  No retained line contains a figure environment, an image-inclusion command or drawing code, so no image is delivered and the package declares no asset dependency.  Editorial formatting only, in addition: an AMS \texttt{amsart} preamble that restates the article's own declarations and macros used by the retained lines, namely the environments \texttt{lemma}, \texttt{theorem} on separate counters, together with the standard packages those lines need.  The retained environments print in this document as Lemma 1, Theorem 1 in order of appearance, whereas the article prints the same items with the numbering of its own full document.  The complete native source of the article as distributed in the arXiv e-print for this exact version is bundled unchanged as \texttt{sources/167970/source0.tex}.
\begin{lemma}\label{lem:isomorphism} $\approx_{\mathsf{CDVS}}$ is Borel reducible to an isomorphism relation of countable structures. Consequently, $\approx_{\mathsf{CDVS}}$ is Borel reducible to $E^\infty_{S_\infty}$.
\end{lemma}

\begin{theorem}\label{thm:CDVSinfty} Let $\Delta$ and $\Lambda$ be countable distance value sets with $\inf\Delta=\inf\Lambda=1$ and $\sup\Delta=\sup\Lambda=\infty$. Then $\Delta$ and $\Lambda$ are equivalent if and only if $\Delta=\Lambda$.
\end{theorem}

\begin{proof} The implication $\Leftarrow$ is obvious. We only show the $\Rightarrow$ direction. For this, let $f$ be a bijection between $\Delta$ and $\Lambda$ witnessing their equivalence. 

For a subset $I\subseteq \mathbb{R}$, let $\Delta_I=\Delta\cap I$ and $\Lambda_I=\Lambda\cap I$. Let $\mathcal{A}_\Delta=(\Delta, R^\Delta)$ and $\mathcal{A}_\Lambda=(\Lambda, R^\Lambda)$ be the countable structures defined in the proof of Lemma~\ref{lem:isomorphism}. Then $f$ is an isomorphism witnessing $\mathcal{A}_\Delta\cong \mathcal{A}_\Lambda$. 

First note that $\inf\Delta=1\in \Delta$ if and only if $\mathcal{A}_\Delta$ satisfies the sentence
$$ \exists x\, \forall y\, [ R(x, x, y)\rightarrow \forall z\, R(z, z, y)]. $$
Similarly for $\Lambda$. Thus $1\in \Delta$ if and only if $1\in \Lambda$. We consider two cases.

{\sc Case 1}: $1\in \Delta\cap \Lambda$. Since $\Delta$ and $\Lambda$ are distance value sets with $\sup\Delta=\inf\Delta=\infty$, we have that $k\in \Delta\cap \Lambda$ for all positive integers $k$. Note that $1$ is definable in $\mathcal{A}_\Delta$ by the formula
$$ \chi_1(x)=\forall y\, [ R(x, x, y)\rightarrow \forall z\, R(z, z, y)]. $$
Similarly for $\mathcal{A}_\Lambda$. Thus $f(1)=1$. Also, we have that $\Delta_{[1,2]}$ is definable in $\mathcal{A}_\Delta$ by the formula
$$ \chi_{[1,2]}(x)=\forall y\, R(y,y,x). $$
Similarly for $\Lambda_{[1,2]}$. Hence $f$ sends $\Delta_{[1,2]}$ to $\Lambda_{[1,2]}$. 

Now for $k\geq 2$, define inductively 
$$ \chi_{[1,k+1]}(x)=\exists y\, [\chi_{[1,k]}(y)\wedge R(1,y,x)]. $$
Then $\Delta_{[1,k+1]}$ is definable in $\mathcal{A}_\Delta$ by $\chi_{[1,k+1]}$, and similarly for $\Lambda_{[1,k+1]}$. It follows that $f$ sends $\Delta_{[1,k+1]}$ to $\Lambda_{[1,k+1]}$ for all $k\geq 1$.

Next we define, by induction on $k\in\mathbb{N}^+$ and $m\in\mathbb{N}$, a formula $\chi_{[1,1+\frac{k}{2^m}]}$ as follows.
$$\begin{array}{rcl} \chi_{[1, 1+\frac{k}{2^0}]}(x)&=&\chi_{[1, k+1]}(x) \mbox{ for all $k\geq 1$,} \\
\chi_{[1,1+\frac{k}{2^{m+1}}]}(x)&=& \forall y\, [R(x,x,y)\rightarrow \chi_{[1,1+\frac{k+2^m}{2^m}]}(y)].
\end{array}
$$
For any $k\in\mathbb{N}^+$ and $m\in\mathbb{N}$, the formula $\chi_{[1, 1+\frac{k}{2^m}]}$ defines $\Delta_{[1,1+\frac{k}{2^m}]}$ in $\mathcal{A}_\Delta$ and $\Lambda_{[1,1+\frac{k}{2^m}]}$ in $\mathcal{A}_\Lambda$. Hence $f$ sends $\Delta_{[1,1+\frac{k}{2^m}]}$ to $\Lambda_{[1,1+\frac{k}{2^m}]}$. This implies that for any dyadic rationals $1<a<b$, $f$ sends $\Delta_{[a,b]}$ to $\Lambda_{[a,b]}$. Hence $f$ must be the identity function, and so $\Delta=\Lambda$.

{\sc Case 2}: $1\not\in \Delta\cup\Lambda$. The formula $\chi_{[1,2]}$ now defines $\Delta_{(1,2]}$ in $\mathcal{A}_\Delta$ and $\Lambda_{(1,2]}$ in $\mathcal{A}_\Lambda$. Thus $f$ sends $\Delta_{(1,2]}$ to $\Lambda_{(1,2]}$. 

Let
$$ \varphi_2(x)=\exists y\, [\chi_{[1,2]}(y)\wedge\forall z\,(\chi_{[1,2]}(z)\rightarrow R(x,y,z))]\wedge \neg \chi_{[1,2]}(x) $$
and let $D_2$ be defined by $\varphi_2$ in $\mathcal{A}_\Delta$. Then $D_2\subseteq\Delta_{(2,3]}$. We claim that $D_2\neq \varnothing$.  To see this, let $a, b\in \Delta_{(1,2]}$ be such that $a<b$ and $2(a-1)<b-1$. Such $a, b$ exist since $\inf\Delta=1\not\in \Delta$.  Then $2a\in\Delta$ since $\Delta$ is a distance value set, and $\varphi_2[2a]$ holds in $\mathcal{A}_\Delta$ as witnessed by $b$. Thus $2a\in D_2$. A similar argument shows $\inf D_2=2$. It follows that $\Delta_{(2,3]}$ is definable in $\mathcal{A}_\Delta$ by the formula
$$ \chi_{(2,3]}(x)=\forall y\, \forall z\, [(\chi_{[1,2]}(y)\wedge \varphi_2(z))\rightarrow R(x,y,z)]\wedge \neg\chi_{[1,2]}(x) $$
and $\Delta_{(1,3]}$ is definable in $\mathcal{A}_\Delta$ by the formula
$$ \chi_{(1,3]}(x)=\chi_{[1,2]}(x)\vee \chi_{(2,3]}(x). $$
Similarly for $\mathcal{A}_\Lambda$. Therefore $f$ sends $\Delta_{(2,3]}$ to $\Lambda_{(2,3]}$ and sends $\Delta_{(1,3]}$ to $\Lambda_{(1,3]}$.

Now for $k\geq 2$, define inductively
$$\begin{array}{rcl} \varphi_{k+1}(x)&=&\exists y\, [\varphi_k(y) \wedge \forall z\, (\chi_{[1,2]}(z)\rightarrow R(x,y,z))]\wedge \neg\chi_{(1,k+1]}(x), \\
\chi_{(k+1,k+2]}(x)&=&\forall y\, \forall z\, [(\chi_{[1,2]}(y)\wedge \varphi_{k+1}(z))\rightarrow R(x,y,z)]\wedge \neg\chi_{(1,k+1]}(x), \\
\chi_{(1,k+2]}(x)&=& \chi_{(1,k+1]}(x)\vee \chi_{(k+1,k+2]}(x). 
\end{array} $$ 
Then for all $k\geq 1$, $\chi_{(1,k+2]}$ defines $\Delta_{(1,k+2]}$ in $\mathcal{A}_\Delta$ and $\Lambda_{(1,k+2]}$ in $\mathcal{A}_\Lambda$. Hence for all $k\geq 2$, $f$ sends $\Delta_{(1,k+1]}$ to $\Lambda_{(1,k+1]}$. 

The rest of the proof is similar to Case 1. We conclude that $f$ must be the identity function, and so $\Delta=\Lambda$.
\end{proof}

\end{document}
